\chapter{Θεωρητικό Υπόβαθρο}
\label{chap:theoretical-background}

\section{Σκοπός και οργάνωση του κεφαλαίου}

Το κεφάλαιο αυτό παρουσιάζει τις έννοιες που είναι απαραίτητες για την κατανόηση
της μεθοδολογίας. Δεν λειτουργεί ως γενική επισκόπηση της ενεργειακής πολιτικής
ή της τεχνητής νοημοσύνης. Εστιάζει στα σημεία που συνδέονται άμεσα με το
αντικείμενο της διπλωματικής: ενεργειακή αναβάθμιση κτιρίων, χρηματοδοτικά
εργαλεία, κατακερματισμένη πληροφορία, δομημένη εξαγωγή δεδομένων και μεγάλα
γλωσσικά μοντέλα.

Η σχέση των επιμέρους εννοιών αποτυπώνεται στο
Σχήμα~\ref{fig:theoretical-stack}. Η λογική του κεφαλαίου είναι διαδοχική:
πρώτα παρουσιάζεται το ενεργειακό και χρηματοδοτικό πλαίσιο, στη συνέχεια η
ανάγκη τυποποίησης της πληροφορίας και, τέλος, οι τεχνικές τεχνητής νοημοσύνης
που μπορούν να υποστηρίξουν αυτή τη μετάβαση.

\begin{figure}[!htbp]
    \centering
    \begin{adjustbox}{max width=\textwidth,center}
        \begin{tikzpicture}[
    layer/.style={
        rectangle,
        rounded corners=3pt,
        draw=blue!52,
        fill=blue!6,
        align=center,
        text width=11.25cm,
        minimum height=0.72cm,
        font=\sffamily\small
    },
    note/.style={
        rectangle,
        rounded corners=2pt,
        draw=gray!35,
        fill=gray!6,
        align=center,
        text width=2.45cm,
        minimum height=0.48cm,
        font=\sffamily\scriptsize
    },
    arrow/.style={-{Latex[length=1.8mm]}, thick, draw=gray!55}
]
\node[layer] (l1) at (0,0) {Ενεργειακή απόδοση κτιρίων και ανάγκη ανακαινίσεων};
\node[layer, below=0.24cm of l1] (l2) {Χρηματοδοτικά εργαλεία και εμπόδιο αρχικού κεφαλαίου};
\node[layer, below=0.24cm of l2] (l3) {Κατακερματισμένες διαδικτυακές πηγές και ετερογενής ορολογία};
\node[layer, below=0.24cm of l3] (l4) {Συλλογή περιεχομένου, καθαρισμός και τυποποίηση};
\node[layer, below=0.24cm of l4] (l5) {LLMs για ταξινόμηση συνάφειας και εξαγωγή δομημένων πεδίων};
\node[layer, below=0.24cm of l5] (l6) {Έλεγχος ποιότητας, τεκμηρίωση πηγής και ανθρώπινη εποπτεία};

\draw[arrow] (l1) -- (l2);
\draw[arrow] (l2) -- (l3);
\draw[arrow] (l3) -- (l4);
\draw[arrow] (l4) -- (l5);
\draw[arrow] (l5) -- (l6);

\node[note, below=0.48cm of l6, xshift=-4.35cm] (n1) {Επίσημες πηγές};
\node[note, below=0.48cm of l6, xshift=-1.45cm] (n2) {Δομημένο JSON};
\node[note, below=0.48cm of l6, xshift=1.45cm] (n3) {Αποθήκευση JSON};
\node[note, below=0.48cm of l6, xshift=4.35cm] (n4) {Web εφαρμογή};
\end{tikzpicture}

    \end{adjustbox}
    \caption[Θεωρητική στοίβα της διπλωματικής εργασίας]{Θεωρητική στοίβα της διπλωματικής εργασίας: από το πρόβλημα ενεργειακής αναβάθμισης και χρηματοδότησης έως την τεχνική επεξεργασία, την τυποποίηση και την παρουσίαση των δεδομένων. Ιδία επεξεργασία.}
    \label{fig:theoretical-stack}
\end{figure}

\section{Ενεργειακή αναβάθμιση και χρηματοδοτικό πλαίσιο}

Η ενεργειακή απόδοση ενός κτιρίου αφορά την ενέργεια που απαιτείται για
θέρμανση, ψύξη, ζεστό νερό χρήσης, φωτισμό και λοιπές λειτουργίες σε σχέση με
τις υπηρεσίες που παρέχονται στους χρήστες. Εξαρτάται από το κτιριακό κέλυφος,
τα τεχνικά συστήματα, τη χρήση του χώρου και τις κλιματικές συνθήκες. Η
διεθνής και ευρωπαϊκή τεκμηρίωση αντιμετωπίζει τον κτιριακό τομέα ως κρίσιμο
πεδίο εξοικονόμησης ενέργειας και μείωσης εκπομπών
\cite{iea_buildings_2026,ipcc_ar6_wg3_buildings,ec_epbd_page_2026,
ec_epbd_news_2025,eea_building_emissions_2025}.

Το πολιτικό πλαίσιο της Ευρωπαϊκής Ένωσης συνδέει την αναβάθμιση του
κτιριακού αποθέματος με την κλιματική ουδετερότητα, την εξοικονόμηση ενέργειας
και την ανάγκη αύξησης του ρυθμού ανακαινίσεων. Σε αυτό εντάσσονται η
Ευρωπαϊκή Πράσινη Συμφωνία, η αναθεωρημένη Οδηγία για την Ενεργειακή Απόδοση
των Κτιρίων, η Οδηγία για την Ενεργειακή Απόδοση, τα εθνικά σχέδια ανακαίνισης
κτιρίων και η στρατηγική \eng{Renovation Wave}
\cite{ec_green_deal_2019,eurlex_epbd_2024,eurlex_eed_2023,ec_nbrp_2026,
ec_renovation_wave_2026,ec_renovation_wave_com_2020}. Επειδή όμως οι
ανακαινίσεις απαιτούν αρχικό κεφάλαιο, η χρηματοδότηση αποτελεί μέρος του
ίδιου προβλήματος και όχι εξωτερική παράμετρο
\cite{ec_financing_renovations_2026,ec_energy_efficiency_financing_report_2026}.

Στην ελληνική περίπτωση, η τεκμηρίωση του προβλήματος και των διαθέσιμων
εργαλείων προκύπτει από στατιστικές πηγές, θεσμικά κείμενα και επίσημες σελίδες
προγραμμάτων. Η ΕΛΣΤΑΤ και η Eurostat περιγράφουν στοιχεία για το κτιριακό
απόθεμα και τις συνθήκες διαβίωσης
\cite{elstat_buildings_census_2021,eurostat_energy_2026,
eurostat_living_conditions_energy_efficiency_2026}, ενώ το ΥΠΕΝ τεκμηριώνει το
πλαίσιο ΚΕΝΑΚ/ΠΕΑ και τη μακροπρόθεσμη στρατηγική ανακαίνισης
\cite{ypen_kenak_2026,ypen_ltrs_2026}. Τα προγράμματα «Εξοικονομώ», οι δράσεις
του Ταμείου Ανάκαμψης και τα τραπεζικά προϊόντα πράσινης χρηματοδότησης
αποτελούν τις βασικές κατηγορίες πηγών που τροφοδοτούν την παρούσα εργασία
\cite{exoikonomo2025_home,exoikonomo2025_programme,greece20_exoikonomo2025,
eurobank_exoikonomo2025_2026,nbg_green_loan_2026}.

Οι κύριες κατηγορίες χρηματοδοτικών εργαλείων και τα πεδία που χρειάζεται να
τυποποιηθούν συνοψίζονται στον Πίνακα~\ref{tab:financing-taxonomy}.

\begin{table}[!htbp]
    \centering
    \caption[Κατηγορίες χρηματοδοτικών εργαλείων και πεδία τυποποίησης]{Βασικές κατηγορίες χρηματοδοτικών εργαλείων και πληροφορία που χρειάζεται τυποποίηση.}
    \label{tab:financing-taxonomy}
    \begin{tabularx}{\textwidth}{L{3.2cm}Y Y}
        \toprule
        \textbf{Κατηγορία} & \textbf{Ενδεικτικό περιεχόμενο} & \textbf{Πεδία ενδιαφέροντος} \\
        \midrule
        Δημόσια προγράμματα & Επιχορηγήσεις, ειδικές δράσεις, προγράμματα τύπου «Εξοικονομώ». & Δικαιούχοι, ποσοστά επιδότησης, προθεσμίες, επιλέξιμες παρεμβάσεις, ενεργειακός στόχος. \\
        Τραπεζικά προϊόντα & Πράσινα ή επισκευαστικά δάνεια με ενεργειακό αντικείμενο. & Επιτόκιο, διάρκεια, ποσό δανείου, σύνδεση με πρόγραμμα, απαιτούμενα δικαιολογητικά. \\
        Ιδιωτικά σχήματα & Συμβάσεις ενεργειακής απόδοσης ή σχήματα ενεργειακών υπηρεσιών. & Μοντέλο σύμβασης, αποπληρωμή, μέτρηση εξοικονόμησης, εμπλεκόμενοι φορείς. \\
        Θεσμικά εργαλεία & Ευρωπαϊκά ή εθνικά ταμεία και πλαίσια πολιτικής. & Φορέας διαχείρισης, επιλέξιμες δράσεις, προϋπολογισμός, γεωγραφικό πεδίο. \\
        \bottomrule
    \end{tabularx}
\end{table}

\section{Από την κατακερματισμένη πληροφορία στα δομημένα δεδομένα}

Η πληροφορία για χρηματοδοτικά εργαλεία δεν εμφανίζεται σε ενιαία μορφή.
Μπορεί να βρίσκεται σε επίσημες ιστοσελίδες, οδηγούς εφαρμογής, PDF,
τραπεζικές σελίδες, δελτία τύπου ή ανακοινώσεις. Σε ευρωπαϊκό επίπεδο, το EU
Building Stock Observatory δείχνει την αξία της οργανωμένης πληροφορίας για
το κτιριακό απόθεμα, την κατανάλωση ενέργειας, τα πιστοποιητικά και τη
χρηματοδότηση \cite{ec_bso_2026}. Η διπλωματική κινείται σε πιο επιχειρησιακή
κλίμακα: επιδιώκει να οργανώσει πληροφορία για συγκεκριμένα ελληνικά
προγράμματα και προϊόντα.

Η διαδικτυακή συλλογή δεδομένων είναι το πρώτο τεχνικό στάδιο αυτής της
μετάβασης. Εργαλεία όπως το Beautiful Soup επιτρέπουν την ανάκτηση και
ανάλυση HTML \cite{beautifulsoup_docs}, ενώ πλαίσια όπως το Scrapy
χρησιμοποιούνται σε πιο σύνθετες ροές crawling \cite{scrapy_docs}. Η συλλογή
πρέπει να γίνεται με σεβασμό στους κανόνες πρόσβασης και στον φόρτο των
εξυπηρετητών. Το Robots Exclusion Protocol αποτελεί σχετική τεχνική αναφορά για
τη συμπεριφορά αυτοματοποιημένων crawlers \cite{rfc9309_robots}.

Η τυποποίηση μετατρέπει το συλλεγμένο κείμενο σε κοινή αναπαράσταση. Στην
εργασία αυτή η αναπαράσταση υλοποιείται ως JSON, επειδή είναι αναγνώσιμο,
διαλειτουργικό και κατάλληλο για διαδικτυακές εφαρμογές. Η λογική της
επικύρωσης σχήματος (\eng{schema validation}) είναι σημαντική: η έξοδος ενός
μοντέλου πρέπει να ελέγχεται απέναντι σε ρητές προσδοκίες δομής, τύπων και
πεδίων \cite{json_schema_spec_2026}. Σε πιο
σύνθετα συστήματα, αντίστοιχες ανάγκες μπορούν να οδηγήσουν σε γράφους γνώσης,
όπου οι οντότητες και οι σχέσεις τους αναπαρίστανται ρητά
\cite{hogan2021knowledge}. Για το συγκεκριμένο σύστημα, η απλούστερη JSON
αναπαράσταση επαρκεί για την αποθήκευση, τον έλεγχο και την παρουσίαση των
προγραμμάτων.

Η ποιότητα των δεδομένων δεν ταυτίζεται μόνο με την ακρίβεια. Περιλαμβάνει
πληρότητα, συνέπεια, επικαιρότητα και καταλληλότητα για χρήση
\cite{wang1996dataquality,batini2009dataquality}. Για το συγκεκριμένο πεδίο,
αυτό σημαίνει ότι ένα πεδίο που δεν υπάρχει στην πηγή πρέπει να παραμένει κενό
ή να επισημαίνεται ως άγνωστο, όχι να συμπληρώνεται με εικασία. Η διάκριση αυτή
είναι κρίσιμη, επειδή η πληροφορία αφορά όρους συμμετοχής, επιτόκια, ποσά και
προθεσμίες.

\section{Τεχνητή νοημοσύνη και μεγάλα γλωσσικά μοντέλα}

Η τεχνητή νοημοσύνη μελετά συστήματα που αντιλαμβάνονται, συλλογίζονται και
δρουν με βάση στόχους, ενώ η μηχανική μάθηση εστιάζει σε μοντέλα που
βελτιώνουν τη συμπεριφορά τους μέσω δεδομένων \cite{russell2021aima,
bishop2006pattern,alpaydin2020ml}. Η βαθιά μάθηση χρησιμοποιεί πολυεπίπεδα
νευρωνικά δίκτυα και έχει αποτελέσει βασικό υπόβαθρο για σύγχρονες εφαρμογές
επεξεργασίας φυσικής γλώσσας \cite{goodfellow2016deep,aggarwal2018neural,
lecun2015deep}. Η επεξεργασία φυσικής γλώσσας παρέχει τις έννοιες για
ταξινόμηση κειμένων, εξαγωγή πληροφορίας, περίληψη και ερωταπαντήσεις
\cite{jurafsky2025slp}.

Η αρχιτεκτονική Transformer εισήγαγε μηχανισμούς προσοχής που επιτρέπουν στα
μοντέλα να συσχετίζουν διαφορετικά σημεία μιας ακολουθίας κειμένου
\cite{vaswani2017attention}. Με βάση αυτή τη λογική, μοντέλα όπως το BERT
ενίσχυσαν την κατανόηση κειμένου \cite{devlin2019bert}, ενώ μοντέλα της
οικογένειας GPT ανέδειξαν την ικανότητα εκτέλεσης εργασιών με λίγα παραδείγματα
\cite{brown2020language}. Για τη διπλωματική, η σημασία τους δεν βρίσκεται στην
ελεύθερη παραγωγή κειμένου, αλλά στην ικανότητά τους να αντιστοιχίζουν
διαφορετικές διατυπώσεις στο ίδιο εννοιολογικό πεδίο.

Η διάκριση αυτή είναι ουσιώδης. Στην ελεύθερη παραγωγή κειμένου, το μοντέλο
μπορεί να συνθέσει απάντηση με βάση τη γενική παραμετρική του γνώση. Στην
παρούσα εργασία ζητείται κάτι διαφορετικό: εξαγωγή τεκμηριωμένη στην πηγή
(\eng{source-backed} ή \eng{source-grounded}). Επομένως, το ζητούμενο δεν είναι να παραχθεί πειστικό
κείμενο, αλλά να μετατραπεί διαθέσιμο περιεχόμενο σε ελεγχόμενα πεδία, με
δυνατότητα επιστροφής στο αρχικό URL.

Οι βασικές τεχνικές αξιοποίησης LLMs που σχετίζονται με την εργασία
συνοψίζονται στον Πίνακα~\ref{tab:llm-techniques}. Στην υλοποίηση δίνεται
έμφαση στον σχεδιασμό προτροπών, στη δομημένη έξοδο και στη σύνδεση της
απάντησης με την πηγή, ενώ τεχνικές όπως το fine-tuning ή η αποδοτική
προσαρμογή παραμέτρων παραμένουν χρήσιμες ως πιθανές μελλοντικές επεκτάσεις.

\begin{table}[!htbp]
    \centering
    \caption[Τεχνικές αξιοποίησης LLMs και σχέση με την εργασία]{Τεχνικές αξιοποίησης LLMs και σχέση τους με την παρούσα εργασία.}
    \label{tab:llm-techniques}
    \begin{tabularx}{\textwidth}{L{3.1cm}Y Y}
        \toprule
        \textbf{Τεχνική} & \textbf{Βασική ιδέα} & \textbf{Σχέση με τη διπλωματική} \\
        \midrule
        Σχεδιασμός προτροπών & Σαφείς οδηγίες και μορφή εξόδου προς το μοντέλο \cite{liu2023prompt}. & Χρησιμοποιείται για συνάφεια, εξαγωγή πεδίων και ερωταπαντήσεις πάνω σε συγκεκριμένα δεδομένα. \\
        RAG και σύνδεση με πηγή & Σύνδεση γλωσσικού μοντέλου με εξωτερικό περιεχόμενο \cite{lewis2020rag}. & Το μοντέλο εργάζεται πάνω σε συλλεγμένη πηγή και όχι σε ελεύθερη γνώση. \\
        Fine-tuning/PEFT & Προσαρμογή μοντέλου σε ειδικό πεδίο με πλήρη ή αποδοτικότερη εκπαίδευση \cite{hu2021lora,houlsby2019adapters,lester2021prompt}. & Δεν αποτελεί βασική επιλογή της υλοποίησης, αλλά μπορεί να εξεταστεί σε επόμενο στάδιο. \\
        Instruction tuning/RLHF & Εκπαίδευση μοντέλων ώστε να ακολουθούν οδηγίες και ανθρώπινες προτιμήσεις \cite{christiano2017preferences,ouyang2022training,bai2022helpful}. & Εξηγεί γιατί τα σύγχρονα LLMs είναι κατάλληλα για εργασίες οδηγιών, χωρίς να αναιρείται η ανάγκη ελέγχου. \\
        \bottomrule
    \end{tabularx}
\end{table}

\section{Αξιοπιστία, κόστος και ανθρώπινη εποπτεία}

Τα LLMs μπορούν να παράγουν απαντήσεις που είναι γλωσσικά πειστικές αλλά όχι
πλήρως τεκμηριωμένες. Το φαινόμενο των hallucinations αποτελεί βασικό
περιορισμό σε εφαρμογές όπου απαιτείται πιστότητα στην πηγή
\cite{ji2023hallucination}. Στο πεδίο των χρηματοδοτικών εργαλείων, τέτοιο
σφάλμα μπορεί να οδηγήσει σε λάθος επιτόκιο, εσφαλμένη προθεσμία ή λανθασμένη
ερμηνεία δικαιούχων. Για αυτό η ανθρώπινη εποπτεία δεν είναι τυπική προσθήκη,
αλλά απαραίτητο μέρος της ροής.

Η ανάγκη αυτή συμφωνεί με ευρύτερες κατευθύνσεις διαχείρισης κινδύνων για
συστήματα τεχνητής νοημοσύνης, όπως το AI Risk Management Framework του NIST
\cite{nist_ai_rmf_2023}. Στη διπλωματική μεταφράζεται σε απλές αλλά κρίσιμες
σχεδιαστικές επιλογές: διατήρηση URL, χρονική σήμανση, κατάσταση ελέγχου,
δυνατότητα διόρθωσης από διαχειριστή και δημοσίευση μόνο ελεγμένων εγγραφών.

Παράλληλα, η χρήση μεγάλων μοντέλων έχει υπολογιστικό και περιβαλλοντικό
κόστος. Η βιβλιογραφία έχει επισημάνει ότι η εκπαίδευση και χρήση μεγάλων
μοντέλων NLP συνδέεται με κατανάλωση ενέργειας και εκπομπές, οι οποίες
εξαρτώνται από την αρχιτεκτονική, το υλικό και το ενεργειακό μίγμα
\cite{strubell2019energy,patterson2021carbon}. Η επιλογή τοπικής εκτέλεσης
μέσω Ollama επιτρέπει ελεγχόμενο πειραματισμό και επαναληψιμότητα, παρότι δεν
αναιρεί την ανάγκη προσεκτικής χρήσης των μοντέλων \cite{ollama_docs_2026}.

Οι βασικοί κίνδυνοι και οι μηχανισμοί μετριασμού συνοψίζονται στον
Πίνακα~\ref{tab:llm-risks}.

\begin{table}[!htbp]
    \centering
    \caption[Κίνδυνοι χρήσης LLMs και μηχανισμοί μετριασμού]{Κίνδυνοι χρήσης LLMs σε χρηματοδοτική πληροφορία και μηχανισμοί μετριασμού.}
    \label{tab:llm-risks}
    \begin{tabularx}{\textwidth}{L{3.3cm}Y Y}
        \toprule
        \textbf{Κίνδυνος} & \textbf{Πιθανή συνέπεια} & \textbf{Μηχανισμός μετριασμού} \\
        \midrule
        Hallucination & Πεδίο που δεν προκύπτει από την πηγή. & Σύνδεση με την πηγή, κενά πεδία όταν η πληροφορία δεν υπάρχει, ανθρώπινος έλεγχος. \\
        Παρωχημένη πληροφορία & Χρήση παλιάς προθεσμίας ή οδηγού. & Χρονική σήμανση και δυνατότητα επανεπεξεργασίας URL. \\
        Σύγχυση οικονομικών όρων & Μπέρδεμα επιδότησης, δανείου, προϋπολογισμού ή επιλέξιμης δαπάνης. & Κοινό σχήμα, κανόνες συνέπειας και έλεγχος από διαχειριστή. \\
        Υπολογιστικό κόστος & Περιορισμένη δυνατότητα μαζικής ή συχνής εκτέλεσης. & Προεπιλογή πηγών, φίλτρο συνάφειας και χρήση κατάλληλου μοντέλου ανά στάδιο. \\
        \bottomrule
    \end{tabularx}
\end{table}

\section{Σύνθεση θεωρητικού υποβάθρου}

Το θεωρητικό υπόβαθρο οδηγεί σε μια συγκεκριμένη μεθοδολογική θέση. Η
ενεργειακή αναβάθμιση χρειάζεται χρηματοδοτικά εργαλεία, τα χρηματοδοτικά
εργαλεία χρειάζονται σαφή και ενημερωμένη πληροφορία, και η πληροφορία αυτή
είναι σήμερα κατακερματισμένη σε διαφορετικές πηγές. Τα μεγάλα γλωσσικά
μοντέλα μπορούν να βοηθήσουν στην ανάγνωση και τυποποίηση αυτής της
πληροφορίας, αλλά μόνο όταν περιορίζονται από την πηγή και συνοδεύονται από
έλεγχο.

Με βάση αυτή τη σύνθεση, η διπλωματική δεν αντιμετωπίζει το LLM ως αυτόνομο
σύμβουλο. Το εντάσσει σε ροή που ξεκινά από επίσημες ή αξιόπιστες πηγές,
καθαρίζει το περιεχόμενο, αποφασίζει τη συνάφεια, εξάγει δομημένα πεδία,
διατηρεί μεταδεδομένα και επιτρέπει ανθρώπινη επιβεβαίωση πριν από την
παρουσίαση στον χρήστη. Η μεθοδολογία που ακολουθεί μεταφέρει αυτή τη λογική
σε συγκεκριμένη ροή εργασίας.
